\begin{Solbook}{1}
 \par \noindent  \begin {enumerate} \item Điều kiện xác định: $x\ne 1$, $x\ne -2$. \begin {align*} \dfrac {3x+4}{x-1}-\dfrac {x-2}{x+2}&=\dfrac {19x+14}{x^2+x-2}\\ (3x+4)(x+2)-(x-2)(x-1)&=19x+14\\ 2x^2-6x-8&=0\\ (x-4)(x+1)&=0. \end {align*} Suy ra $x=4$ hoặc $x=-1$; cả hai giá trị đều thỏa mãn điều kiện xác định. Vậy $S=\{-1;4\}$. \par \item Nhân cả hai vế với $4>0$, ta được \begin {align*} -8(3x-5)+2(5x+1)&>3-x\\ -14x+42&>3-x\\ 39&>13x\\ x&<3. \end {align*} Vậy bất phương trình có nghiệm $x<3$. \par \item Dùng phương pháp thế, ta có \begin {align*} &\begin {cases} \dfrac {x+2y}{3}-\dfrac {2x-y}{4}=2,\\ 3(x-y)+2(2x+y)=-9 \end {cases}\\[0.7em] \Leftrightarrow \quad &\begin {cases} -2x+11y=24,\\ 7x-y=-9 \end {cases}\\[0.7em] \Leftrightarrow \quad &\begin {cases} -2x+11y=24,\\ y=7x+9 \end {cases}\\[0.7em] \Leftrightarrow \quad &\begin {cases} -2x+11(7x+9)=24,\\ y=7x+9 \end {cases}\\[0.7em] \Leftrightarrow \quad &\begin {cases} 75x=-75,\\ y=7x+9 \end {cases}\\[0.7em] \Leftrightarrow \quad &\begin {cases} x=-1,\\ y=2. \end {cases} \end {align*} Vậy hệ phương trình có nghiệm $(x;y)=(-1;2)$. \end {enumerate}  \hfill \qedEX 
\end{Solbook}
\begin{Solbook}{2}
 \par \noindent  \begin {enumerate} \setlength {\itemsep }{1em} \par \item Ta có \begin {align*} x+3+4\sqrt {x-1}&=\left (\sqrt {x-1}+2\right )^2,\\ x+3-4\sqrt {x-1}&=\left (\sqrt {x-1}-2\right )^2,\\ 4x-4&=\left (2\sqrt {x-1}\right )^2. \end {align*} Vì $1\le x\le 5$ nên $0\le \sqrt {x-1}\le 2$. Do đó \begin {align*} A &=\left |\sqrt {x-1}+2\right | +\left |\sqrt {x-1}-2\right | -\left |2\sqrt {x-1}\right |\\ &=\sqrt {x-1}+2+2-\sqrt {x-1}-2\sqrt {x-1}\\ &=4-2\sqrt {x-1}. \end {align*} Vậy $A=4-2\sqrt {x-1}$. \par \item Với $x\ge 0$ và $x\ne 4$, ta có \begin {align*} B &=\dfrac {(\sqrt {x}-2)(\sqrt {x}+2)}{\sqrt {x}-2} -\dfrac {(\sqrt {x}-1)(\sqrt {x}+1)}{\sqrt {x}+1} +\dfrac {2}{\sqrt {x}-2}-\dfrac {2}{\sqrt {x}+2}\\ &=\sqrt {x}+2-\left (\sqrt {x}-1\right ) +\dfrac {2(\sqrt {x}+2)-2(\sqrt {x}-2)}{x-4}\\ &=3+\dfrac {8}{x-4}\\ &=\dfrac {3(x-4)+8}{x-4}\\ &=\dfrac {3x-4}{x-4}. \end {align*} Vậy $B=\dfrac {3x-4}{x-4}$. \end {enumerate}  \hfill \qedEX 
\end{Solbook}
\begin{Solbook}{3}
 \par \noindent  Gọi $x$ là số học sinh nam và $y$ là số học sinh nữ của lớp, với $x,y$ nguyên dương và $x+y=36$. \par Tổng chiều cao của cả lớp là $160\cdot 36=5760$ (cm). Tổng chiều cao của các bạn nam là $168x$ (cm) và của các bạn nữ là $156y$ (cm), nên \[ 168x+156y=5760. \tag {1} \] Vì $y=36-x$, thay vào $(1)$: \begin {align*} 168x+156(36-x)&=5760\\ 168x+5616-156x&=5760\\ 12x&=144\\ x&=12. \end {align*} Suy ra $y=36-12=24$. \par \emph {Kiểm tra.} Tổng chiều cao là \[ 168\cdot 12+156\cdot 24=2016+3744=5760\ \text {(cm)}, \] trung bình $5760:36=160$ (cm), đúng với giả thiết. \par Vậy lớp có $12$ học sinh nam và $24$ học sinh nữ.  \hfill \qedEX 
\end{Solbook}
\begin{Solbook}{4}
 \par \noindent  \textbf {a)} Bán kính của trái bóng là: \[ 22{,}3:2=11{,}15\text { (cm)}. \] \par Diện tích bề mặt của quả bóng là: \[ S=4\pi \cdot 11{,}15^2 \approx 1562\text { (cm}^2\text {)}. \] \par \textbf {b)} Gọi $x$, $y$ (múi) lần lượt là số múi da màu đen và màu trắng của trái bóng ($x,y\in \mathbb {N}^{*}$). \par Vì trái bóng được may từ $32$ múi da đen và trắng nên ta có phương trình: \[ x+y=32. \qquad (1) \] \par Vì mỗi múi da màu đen có diện tích $37$ cm$^2$ và mỗi múi da màu trắng có diện tích $55{,}9$ cm$^2$ nên theo diện tích bề mặt đã làm tròn ở câu a), ta có phương trình: \[ 37x+55{,}9y=1562. \qquad (2) \] \par Từ (1) và (2), ta có hệ phương trình: \[ \begin {cases} x+y=32\\ 37x+55{,}9y=1562 \end {cases} \] \par \[ \Leftrightarrow \begin {cases} x=32-y\\ 37(32-y)+55{,}9y=1562 \end {cases} \] \par \[ \Leftrightarrow \begin {cases} x=32-y\\ 18{,}9y=378 \end {cases} \Leftrightarrow \begin {cases} x=12\\ y=20 \end {cases} \quad \text {(thỏa mãn).} \] \par Vậy trên trái bóng có $12$ múi da màu đen và $20$ múi da màu trắng.  \hfill \qedEX 
\end{Solbook}
\begin{Solbook}{5}
 \par \noindent  \begin {enumerate} \item Gọi $M$ là vị trí mắt, $C$ là chân cột ăng-ten, $B$ là đỉnh cột và $A$ là chân tường. Kẻ $MH$ vuông góc với $AB$ tại $H$. Khi đó $MH$ song song với mặt đất, $AH=1{,}5$ m và $BC=6$ m. \par Đặt $MH=d$ (m), $d>0$. \par Tam giác $MHC$ vuông tại $H$, $\widehat {HMC}=30^\circ $, nên \[ \tan 30^\circ =\dfrac {HC}{d} \Rightarrow HC=d\cdot \dfrac {\sqrt {3}}{3}. \] Tam giác $MHB$ vuông tại $H$, $\widehat {HMB}=45^\circ $, nên \[ \tan 45^\circ =\dfrac {HB}{d} \Rightarrow HB=d. \] Vì $B$, $C$, $H$ thẳng hàng và $C$ nằm giữa $H$ và $B$ nên $HB=HC+CB$: \begin {align*} d&=d\cdot \dfrac {\sqrt {3}}{3}+6\\ d\left (1-\dfrac {\sqrt {3}}{3}\right )&=6\\ d\cdot \dfrac {3-\sqrt {3}}{3}&=6\\ d(3-\sqrt {3})&=18\\ d&=\dfrac {18}{3-\sqrt {3}}. \end {align*} Nhân liên hợp với $3+\sqrt {3}$: \[ d=\dfrac {18(3+\sqrt {3})}{(3-\sqrt {3})(3+\sqrt {3})} =\dfrac {18(3+\sqrt {3})}{6} =3(3+\sqrt {3}) =9+3\sqrt {3}. \] Vậy Minh đứng cách mặt tường tòa nhà $9+3\sqrt {3}$ m. \par \item Chiều cao nóc nhà là \begin {align*} AC &=AH+HC\\ &=1{,}5+d\cdot \dfrac {\sqrt {3}}{3}\\ &=1{,}5+(9+3\sqrt {3})\cdot \dfrac {\sqrt {3}}{3}\\ &=1{,}5+\dfrac {9\sqrt {3}+9}{3}\\ &=1{,}5+3\sqrt {3}+3\\ &=4{,}5+3\sqrt {3}\ \text {(m)}. \end {align*} Vậy nóc nhà cao $4{,}5+3\sqrt {3}$ m so với mặt đất. \end {enumerate}  \hfill \qedEX 
\end{Solbook}
